% changelog: "0.6.0, 2024-05-28, Guglielmo Barison, Stesura liste di controllo"

\documentclass[8pt]{article}
\usepackage[italian]{babel}
\usepackage[utf8]{inputenc}
\usepackage[letterpaper, left=1in, right=1in, bottom=0.75in, top=0.75in]{geometry}
\usepackage{amsmath}
\usepackage{subfiles}
\usepackage{lipsum}
\usepackage{csquotes}
\usepackage{amsfonts}
\usepackage[sfdefault]{plex-sans}
\usepackage{float}
\usepackage{pifont}
\usepackage{mathabx}
\usepackage[euler]{textgreek}
\usepackage{makecell}
\usepackage{tikz}
\usepackage{wrapfig}
\usepackage{siunitx}
\usepackage{amssymb} 
\usepackage{tabularx}
\usepackage{adjustbox}
\usepackage[document]{ragged2e}
\usepackage{floatflt}
\usepackage[hidelinks]{hyperref}
\usepackage{graphicx}
\usepackage{hyperref}
\setcounter{tocdepth}{4}
\usepackage{caption}
\usepackage{multicol}
\usepackage{tikz}
\setlength\parindent{0pt}
\captionsetup{font=footnotesize}
\usepackage{fancyhdr} 
\usepackage{graphicx}
\usepackage{capt-of}% 
\usepackage{booktabs}
\usepackage{varwidth}
\usepackage{datetime2}
\usepackage{xcolor}
\usepackage{longtable}
\usepackage{array}
\usepackage{ragged2e}
\usepackage{colortbl}
\usepackage{verbatim}

\newcommand{\customtitle}{PIANO DI QUALIFICA}% o ESTERNO

% -- STILE COLONNA CENTRATA PER TABELLE -- %
\newcolumntype{M}[1]{>{\centering\arraybackslash}m{#1}}

% -- STILE INTESTAZIONE -- %
\fancypagestyle{mystyle}{
	\fancyhf{} 
	\fancyhead[R]{\includegraphics[height=1cm]{../../template/images/logos/NaN1fy_logo.png}} 
	\fancyhead[L]{\leftmark} 
	\renewcommand{\headrulewidth}{1pt} 
	\fancyhead[L]{\customtitle} 
	\renewcommand{\headsep}{1.3cm} 
	\fancyfoot[C]{\thepage} 
}

% -- PER LA FIRMA -- %
\newcommand{\signatureline}[1]{%
	 \par\vspace{0.5cm}
	\noindent\makebox[\linewidth][r]{\rule{0.2\textwidth}{0.5pt}\hspace{3cm}\makebox[0pt][r]{\vspace{3pt}\footnotesize #1}}%
}

% -- PER IL GLOSSARIO -- %
\newcommand{\glossterm}[1]{#1\textsuperscript{G}} % inserisci \glossterm{termine}

% -- per abilitare 4x sottosezioni es 2.1.1.1
\setcounter{secnumdepth}{4}
\newcommand{\subsubsubsection}[1]{\paragraph{#1}\mbox{}\\\\}

\begin{document}
\definecolor{myblue}{RGB}{23,103,162}
\begin{titlepage}
	\begin{tikzpicture}[remember picture, overlay]
		\node[anchor=south east, opacity=0.2, yshift = -4cm, xshift= 2em] at (current page.south east) {\includegraphics[width=0.7\textwidth, trim=0cm 0cm 5cm 0cm, clip]{../../template/images/logos/Universita_Padova_transparent.png}}; 
		\node[anchor=north west, opacity=1, yshift = 4.2cm, xshift= 1.4cm, scale=1.6] at (current page.south west) {\includegraphics[width=4cm]{../../template/images/logos/NaN1fy_logo.png}};
	\end{tikzpicture}
	
	\begin{minipage}[t]{0.47\textwidth}
		{\large{\textsc{Destinatari}}
			\vspace{3mm}
			\\ \large{\textsc{Prof. Tullio Vardanega}}
			\\ \large{\textsc{Prof. Riccardo Cardin}}
		}
	\end{minipage}
	\hfill
	\begin{minipage}[t]{0.47\textwidth}\raggedleft
		{\large{\textsc{Redattori}}
			\vspace{3mm}
			{\\\large{\textsc{Guglielmo Barison}\\}} % massimo due 
			{\large{\textsc{Davide Donanzan}}}
			
			
		}
		\vspace{8mm}
		
		{\large{\textsc{Verificatori}}
			\vspace{3mm}
			{\\\large{\textsc{XXXX XXXX}\\}} 
			{\large{\textsc{XXXX XXXX}}}
			
		}
		\vspace{4mm}\vspace{4mm}
	\end{minipage}
	\vspace{4cm}
	\begin{center}
		\begin{flushright}
			{\fontsize{30pt}{52pt}\selectfont \textbf{Piano Di Qualifica}} % o ESTERNO
		\end{flushright}
		\vspace{3cm}
	\end{center}
	\vspace{10 cm}
	{\small \textsc{\href{mailto: nan1fyteam.unipd@gmail.com}{nan1fyteam.unipd@gmail.com}}}
\end{titlepage}
\pagestyle{mystyle}
\section*{Registro delle Modifiche}
\begin{table}[ht!]	
	\centering
	\begin{tabular}{p{1.2cm} p{2cm} p{6cm} p{3cm} p{2cm}}
		\toprule
		\textbf{Versione}& \textbf{Data} & \textbf{Descrizione} & \textbf{Autore} & \textbf{Ruolo} \\
		\midrule	
        0.6.0 & 2024-05-28 & Stesura sezione ``Liste di controllo". & Guglielmo Barison & Redattore
        \\\\
        0.5.0 & 2024-05-24 & Stesura sezione ``Cruscotto della qualità". & Davide Donanzan & Redattore \\\\
        0.4.0 & 2024-05-21 & Stesura sezione ``Test di accettazione". & Guglielmo Barison & Redattore \\\\
        0.3.1 & 2024-05-15 & Modifiche minori alla struttura di ``Test di Unità". & Davide Donanzan & Redattore \\\\
		0.3.0 & 2024-05-14 & Inizio stesura ``Test di Unità" e ``Test di Integrazione". & Davide Donanzan & Redattore \\\\
		0.2.0 & 2024-04-28 & Inizio stesura ``Strategie di testing". & Guglielmo Barison & Redattore \\\\
		0.1.0 & 2024-04-18 & Stesura ``Obiettivi di qualità". & Guglielmo Barison & Redattore \\\\
		0.0.0 & 2024-04-09 & Struttura di base ed introduzione.  & Guglielmo Barison & Redattore \\
		\bottomrule
		% Ruolo Redattore o Verificatore
	\end{tabular}
	\caption*{Tabella: Registro delle modifiche.}
	\label{table:Registro delle modifiche}
\end{table}
\newpage
\tableofcontents
\newpage
\listoffigures
\newpage
\listoftables
\newpage
\justifying
\section{Introduzione}
\subsection{Scopo del documento}
Questo documento offre una panoramica dettagliata delle strategie di verifica e validazione adottate per assicurare la \glossterm{qualità} del prodotto e dei processi nel contesto del progetto in questione. Sarà costantemente aggiornato per riflettere l'evoluzione del progetto e concentrerà l'attenzione sui risultati delle verifiche per risolvere tempestivamente eventuali criticità.
\\\\
Il \textit{Piano di Qualifica}, dinamico e incrementale, illustra le pratiche per il controllo di qualità degli artefatti e dei processi, con particolare enfasi sulle metriche di valutazione del prodotto. È progettato per guidare l'adozione di processi mirati al miglioramento continuo, fornendo misure quantitative per valutare il progresso del progetto. Questo impegno costante per la qualità si riflette nel regolare aggiornamento del documento per adattarsi alle esigenze mutevoli del progetto, garantendo così la crescita e l'evoluzione sia del processo che del prodotto nel tempo.
\subsection{Scopo del capitolato}
Lo scopo del \glossterm{capitolato} C6 è proporre una soluzione per la creazione di una piattaforma di monitoraggio per una smart city. Tale piattaforma deve raccogliere e analizzare in tempo reale una vasta gamma di dati provenienti da sensori distribuiti nella città, riguardanti aspetti quali traffico, qualità dell'aria, consumi energetici e altro ancora. L'obiettivo è fornire alle autorità locali informazioni dettagliate per prendere decisioni informate sulla gestione delle risorse e l'implementazione dei servizi, coinvolgendo anche i cittadini attraverso la condivisione di dati e la partecipazione attiva. La soluzione proposta prevede l'utilizzo di tecnologie per il data streaming processing e la simulazione dei dati dei sensori, con l'obiettivo di fornire dashboard intuitive per la visualizzazione dei dati raccolti e l'analisi delle condizioni della città.\\
L’applicativo sviluppato richiede una copertura dei \glossterm{test} di almeno l’80\%.

\subsection{Glossario}
Per garantire chiarezza nel linguaggio utilizzato nei documenti, è stato redatto un Glossario contenente le definizioni dei termini con significato specifico da disambiguare. Tali termini sono contrassegnati con una G ad apice. L'inserimento di un termine nel Glossario è considerato completo solo dopo averne fornito la definizione.
\subsection{Riferimenti}
\subsubsection{Normativi}
\begin{itemize}
	\item \textit{Norme di progetto v.X.X.X};
	\item Presentazione e documentazione del \glossterm{capitolato} d’appalto C6 - SyncCity:
	\begin{itemize}
		\item \href{https://www.math.unipd.it/~tullio/IS-1/2023/Progetto/C6p.pdf}{\color{myblue}https://www.math.unipd.it\textasciitilde{}tullio/IS-1/2023/Progetto/C6p.pdf} (Ultimo accesso: \today)
		\item \href{https://www.math.unipd.it/~tullio/IS-1/2023/Progetto/C6.pdf}{\color{myblue}https://www.math.unipd.it/\textasciitilde{}tullio/IS-1/2023/Progetto/C6.pdf} (Ultimo accesso: \today)
	\end{itemize}
	\item Regolamento di progetto:
	\begin{itemize}
		\item \href{https://www.math.unipd.it/~tullio/IS-1/2023/Dispense/PD2.pdf}{\color{myblue}https://www.math.unipd.it/\textasciitilde{}tullio/IS-1/2023/Dispense/PD2.pdf} (Ultimo accesso: \today)
	\end{itemize}
\end{itemize}
\clearpage
\subsubsection{Informativi}
\begin{itemize}
	\item Dispense T7 - Qualità del software:
	\begin{itemize}
		\item \href{https://www.math.unipd.it/~tullio/IS-1/2023/Dispense/T7.pdf}{\color{myblue}https://www.math.unipd.it/\textasciitilde{}tullio/IS-1/2023/Dispense/T7.pdf} (Ultimo accesso: \today)
	\end{itemize}
	\item Dispense T8 - Qualità di processo:
	\begin{itemize}
		\item \href{https://www.math.unipd.it/~tullio/IS-1/2023/Dispense/T8.pdf}{\color{myblue}https://www.math.unipd.it/\textasciitilde{}tullio/IS-1/2023/Dispense/T8.pdf} (Ultimo accesso: \today)
	\end{itemize}
	\item Dispense T9 - Verifica e validazione:
	\begin{itemize}
		\item \href{https://www.math.unipd.it/~tullio/IS-1/2023/Dispense/T9.pdf}{\color{myblue}https://www.math.unipd.it/\textasciitilde{}tullio/IS-1/2023/Dispense/T9.pdf} (Ultimo accesso: \today)
	\end{itemize}
	\item \glossterm{ISO} / \glossterm{IEC} 9126:
	\begin{itemize}
		\item \href{https://it.wikipedia.org/wiki/ISO/IEC_9126}{\color{myblue}https://it.wikipedia.org/wiki/ISO/IEC\_9126} (Ultimo accesso: \today)
	\end{itemize}
	\item \glossterm{ISO} / \glossterm{IEC} 12207-1995:
	\begin{itemize}
		\item \href{https://www.math.unipd.it/~tullio/IS-1/2009/Approfondimenti/ISO_12207-1995.pdf}{\color{myblue}https://www.math.unipd.it/\textasciitilde{}tullio/IS-1/2009/Approfondimenti/ISO\_12207-1995.pdf} \\ (Ultimo accesso: \today)
	\end{itemize}
\end{itemize}
\clearpage
\section{Obiettivi di qualità}
Ogni \glossterm{processo} viene valutato tramite l'utilizzo di metriche specifiche, le cui definizioni sono dettagliate nelle sezioni Metriche di qualità del processo e Metriche di qualità del prodotto del documento \textit{Norme di Progetto vX.X.X}. Queste sezioni delineano i criteri che le metriche devono rispettare per essere valutate come accettabili o eccellenti. La sigla MPC sta ad indicare le metriche di processo.
% ricordarsi di verficare effettivamente il nome della sezione in NdP, nome in corsivo per riferimenti nel testo a documentazione del gruppo?
\subsection{Qualità di processo}
La base della \glossterm{qualità} del processo risiede nell'idea che, per ottenere un prodotto conforme a determinati standard di qualità, sia fondamentale sottoporre i processi che lo supportano a controlli regolari, al fine di ottimizzarli. Il concetto di qualità del processo viene quindi applicato a tutte le attività, pratiche e metodologie utilizzate lungo l'intero ciclo di vita del software. In breve, la qualità del processo mira a integrare la qualità nel prodotto stesso, garantendo che sia intrinseca al processo e non solo un obiettivo secondario.
\subsubsection{Processi primari}
\subsubsubsection{Fornitura} 
\begin{table}[h]	
	\centering
	\begin{tabular}{lccc}
		\toprule
		\textbf{Metrica}& \textbf{Descrizione} & \textbf{Valore accettabile} & \textbf{Valore ottimo} \\
		\midrule
		MCP-EV & Earned value (EV) & $\geq$ 0 & $\leq$ EAC \\\\
		MPC-PV & Planned Value (PV) & $\geq$ 0 & $\leq$ BAC\\\\
		MPC-AC & Actual costo (AC) & $\geq$ 0 & $\leq$ EAC\\\\
		MPC-CPI & Cost Performance Index (CPI) & tra 0.95 e 1.05 & $\leq$ 1\\\\
		MPC-EAC & Estimate At Completion (EAC) & deviazione del $\pm$ 5\% dal BAC & BAC\\\\
		MPC-ETC & Estimate To Completion (ETC) & $\geq $ 0 & $\leq$ EAC\\\\
		MPC-VAC & Variance At Completion (VAC) & deviazione del $\pm$ 10\% dal BAC & 0\%\\\\
		MPC-SV & Schedule Variance (SV) & deviazione del $\pm$ 10\% dal BAC & 0\%\\\\
		MPC-BV & Budget Variance (BV) & deviazione del $\pm$ 10\% dal BAC  & 0\%\\
		\bottomrule
		% Ruolo Redattore o Verificatore
	\end{tabular}
	\caption{Metriche per il processo di fornitura.}
	\label{table:Tabella metriche per il processo di fornitura.}
\end{table}
\clearpage
\subsubsubsection{Sviluppo}
\begin{table}[h]	
	\centering
	\begin{tabular}{lccc}
		\toprule
		\textbf{Metrica}& \textbf{Descrizione} & \textbf{Valore accettabile} & \textbf{Valore ottimo} \\
		\midrule
		MCP-RSI & Requirements stability index (RSI) & $\geq $ 75\%  & $\leq$ 100\% \\\\
		MPC-SFIN & Structural Fan-In (SFIN) & - & Va massimizzato\\\\
		MPC-SFOUT & Structural Fan-Out (SFOUT) & - & Va minimizzato\\
		\bottomrule
		% Ruolo Redattore o Verificatore
	\end{tabular}
	\caption{Valori accettabili e ottimi per ogni metrica riguardante il processo di sviluppo.}
	\label{table:Valori accettabili e ottimi per ogni metrica riguardante il processo di sviluppo.}
\end{table}
\subsubsection{Processi di supporto}
\subsubsubsection{Documentazione}
\begin{table}[h]	
	\centering
	\begin{tabular}{lccc}
		\toprule
		\textbf{Metrica}& \textbf{Descrizione} & \textbf{Valore accettabile} & \textbf{Valore ottimo} \\
		\midrule
		MPC-IG & Indice Gulpease (IGI) & $\geq$ 60\% & 100 \\\\
		MPC-CO & Correttezza Ortografica (CO) & 0 & 0 \\
		\bottomrule
		% Ruolo Redattore o Verificatore
	\end{tabular}
	\caption{Metriche per il processo di documentazione.}
	\label{table:Tabella delle metriche per il processo di documentazione}
\end{table}
\subsubsubsection{Verifica}
\begin{table}[h]	
	\centering
	\begin{tabular}{lccc}
		\toprule
		\textbf{Metrica}& \textbf{Descrizione} & \textbf{Valore accettabile} & \textbf{Valore ottimo} \\
		\midrule
		MPC-CC & Code Coverage (CC) & $\geq$ 80\% & 100\% \\\\
		MPC-PT & Passed Test cases percentage (PT) & 100\% & 100\% \\
		\bottomrule
	\end{tabular}
	\caption{Metriche per il processo di documentazione.}
	\label{table:Tabella delle metriche per il processo di documentazione}
\end{table}
\subsubsubsection{Gestione della qualità}
\begin{table}[H]	
	\centering
	\begin{tabular}{lccc}
		\toprule
		\textbf{Metrica}& \textbf{Descrizione} & \textbf{Valore accettabile} & \textbf{Valore ottimo} \\
		\midrule
		MPC-QMS & Quality Metrics Satisfied (QMS) & $\geq$ 85\%& 100\%\\
		\bottomrule
	\end{tabular}
	\caption{Valori accettabili e ottimi per ogni metrica riguardante il processo di gestione della qualità.}
	\label{table:Valori accettabili e ottimi per ogni metrica riguardante il processo di gestione della qualità.}
\end{table}
\clearpage
\subsubsection{Processi organizzativi}
\subsubsubsection{Gestione dei processi}
\begin{table}[H]	
	\centering
	\begin{tabular}{lccc}
		\toprule
		\textbf{Metrica}& \textbf{Descrizione} & \textbf{Valore accettabile} & \textbf{Valore ottimo} \\
		\midrule
		MPC-NR & Non-calculated Risk (NR) & $\leq$ 3 & 0\\\\
		MPC-ET & Efficienza temporale (ET) & $\leq$ 3 & $\leq$ 1 \\
		\bottomrule
	\end{tabular}
	\caption{Valori accettabili e ottimi per ogni metrica riguardante il processo di gestione dei processi.}
	\label{table:Valori accettabili e ottimi per ogni metrica riguardante il processo di gestione dei processi.}
\end{table}
\subsection{Qualità di prodotto}
Si riferisce alle caratteristiche di un'entità risultante dallo sviluppo software, che influenzano la sua capacità di soddisfare sia le esigenze esplicite che implicite. In altre parole, è quanto il prodotto si adatta alle aspettative del cliente o agli standard predefiniti. Questo implica una valutazione completa del software realizzato, concentrandosi su attributi come usabilità, \glossterm{funzionalità}, affidabilità e manutenibilità, oltre alle prestazioni generali. L'obiettivo è garantire che il software non solo soddisfi le richieste del cliente e funzioni correttamente, ma che lo faccia conformemente ai rigidi standard di \glossterm{qualità} stabiliti. Per raggiungere questo obiettivo, il team si impegna a seguire le metriche di prodotto, indicate con la sigla MPD, come specificato nel documento \textit{Norme di Progetto vX.X.X}.
% ricordarsi di verficare effettivamente il nome della sezione in NdP, nome in corsivo per riferimenti nel testo a documentazione del gruppo?
\subsubsection{Funzionalità}
\begin{table}[H]	
	\centering
	\begin{tabular}{lccc}
		\toprule
		\textbf{Metrica}& \textbf{Descrizione} & \textbf{Valore accettabile} & \textbf{Valore ottimo} \\
		\midrule
		MPD-ROS& Requisiti Obbligatori Soddisfatti (ROS) & 100\% & 100\%\\\\
		MPD-RDS & Requisiti Desiderabili Soddisfatti (RDS) & $\geq$ 0\% & $\geq$ 75\% \\\\
		MPD-ROPS & Requisiti Opzionali Soddisfatti (ROPS) & $\geq$ 0\% & $\geq$ 75\% \\
		\bottomrule
	\end{tabular}
	\caption{Valori accettabili e ottimi per ogni metrica riguardante la funzionalità del prodotto.}
	\label{table:Valori accettabili e ottimi per ogni metrica riguardante la funzionalità del prodotto.}
\end{table}
\subsubsection{Affidabilità}
\begin{table}[H]	
	\centering
	\begin{tabular}{lccc}
		\toprule
		\textbf{Metrica}& \textbf{Descrizione} & \textbf{Valore accettabile} & \textbf{Valore ottimo} \\
		\midrule
		MPD-BC & Branch Coverage (BC) & $\geq$ 80\% & 100\%\\\\
		MPD-SC & Statement Coverage (SC) & $\geq$ 80\% & 100\% \\\\
		MPD-FD & Failure Density (FD) & $\geq$ 80\% & 100\% \\\\
		MPD-PTCP & Passed Test Cases Percentage (PTCP) & 100\%  & 100\% \\
		\bottomrule
	\end{tabular}
	\caption{Valori accettabili e ottimi per ogni metrica riguardante l’affidabilità del prodotto.}
	\label{table:Valori accettabili e ottimi per ogni metrica riguardante l’affidabilità del prodotto.}
\end{table}
\subsubsection{Usabilità}
\begin{table}[H]	
	\centering
	\begin{tabular}{lccc}
		\toprule
		\textbf{Metrica}& \textbf{Descrizione} & \textbf{Valore accettabile} & \textbf{Valore ottimo} \\
		\midrule
		MPD-FU & Facilità di Utilizzo (FU) & $\geq$ 9 click & $\geq$ 5 click \\\\
		MPD-TA & Tempo di Apprendimento (TA) & $\leq$ 15 minuti & $\leq$ 5 minuti \\
		\bottomrule
	\end{tabular}
	\caption{Valori accettabili e ottimi per ogni metrica riguardante l’usabilità del prodotto.}
	\label{table:Valori accettabili e ottimi per ogni metrica riguardante l’usabilità del prodotto.}
\end{table}
\subsubsection{Efficienza}
\begin{table}[H]	
	\centering
	\begin{tabular}{lccc}
		\toprule
		\textbf{Metrica}& \textbf{Descrizione} & \textbf{Valore accettabile} & \textbf{Valore ottimo} \\
		\midrule
		MPD-UR & Utilizzo risorse (UR) & $\geq$ 75\% & 100\% \\
		\bottomrule
	\end{tabular}
	\caption{Valori accettabili e ottimi per ogni metrica riguardante l’efficienza del prodotto.}
	\label{table:Valori accettabili e ottimi per ogni metrica riguardante l’efficienza del prodotto.}
\end{table}
\subsubsection{Manutenibilità}
\begin{table}[H]	
	\centering
	\begin{tabular}{lccc}
		\toprule
		\textbf{Metrica}& \textbf{Descrizione} & \textbf{Valore accettabile} & \textbf{Valore ottimo} \\
		\midrule
		MPD-CC & Complessità Ciclomatica (CC) & 11-20 & 1-10 \\\\
		MPD-CS & Code Smell (CS) & 0 & 0 \\
		\bottomrule
	\end{tabular}
	\caption{Valori accettabili e ottimi per ogni metrica riguardante la manutenibilità del prodotto.}
	\label{table:Valori accettabili e ottimi per ogni metrica riguardante la manutenibilità del prodotto.}
\end{table}
\clearpage
\section{Strategie di testing}
In questa sezione viene esposto il piano di testing che verrà utilizzato per garantire la correttezza finale del prodotto. Come enunciato nel documento \textit{Norme di Progetto v.X.X.X}, il piano segue il modello a \glossterm{V}, il quale associa ad ogni fase di sviluppo una corrispondente tipologia di testing. Tali tipologie sono le seguenti:
\begin{itemize}
	\item \textbf{Test di unità:} si verifica il corretto funzionamento delle unità componenti il \glossterm{sistema}. Un’unità rappresenta un elemento indivisibile e indipendente del \glossterm{sistema}; 
	\item \textbf{Test di integrazione:} si verifica il corretto funzionamento di più unità che cooperano per svolgere uno specifico compito (tali unità devono certamente aver superato i loro test di unità precedentemente);
	\item \textbf{\glossterm{Test di sistema}:} si verifica il corretto funzionamento del \glossterm{sistema} nella sua interezza. I requisiti funzionali obbligatori, di vincolo, di \glossterm{qualità} e di prestazione, precedentemente concordati con la Proponente mediante stipulazione del contratto, devono essere soddisfatti per intero;
	\item \textbf{Test di accettazione:} si verifica il soddisfacimento della Proponente rispetto al prodotto software. Il loro superamento permette di procedere con il rilascio del prodotto.
\end{itemize}
Per le procedure necessarie all’esecuzione di test di unità e di integrazione si rimanda al documento \textit{Norme di Progetto vX.X.X} nella sezione relativa al processo di verifica.
% verificare effettivamente poi norme di progetto
\subsection{Codice dei test}
Ogni test è associato ad un codice univoco definito nel seguente formato:
\begin{center}
	\textbf{T[Tipologia]-[Numero]}
\end{center}
Dove \textbf{Tipologia} indica la tipologia del test: 
\begin{itemize}
    \item \textbf{U:} di unit\`{a};
	\item \textbf{I:} di integrazione;
	\item \textbf{S:} di sistema;
	\item \textbf{A:} di accettazione.
\end{itemize}
Ogni test ha uno \textbf{Stato}, che puo essere:
\begin{itemize}
	\item \textbf{V:} verificato. Il test ha esito positivo;
	\item \textbf{NV:} non verificato. Il test ha esito negativo; 
	\item \textbf{NI:} non implementato.
\end{itemize}
\clearpage
\subsection{Test di unità}
\renewcommand{\arraystretch}{2.5}
\rowcolors{2}{gray!20}{white}
\begin{longtable}{|>{\centering}p{2cm}|>{\RaggedRight}m{12cm}|>{\centering\arraybackslash}p{2cm}|}
    \hline
    \rowcolor{white}
    \textbf{Codice Test} & \textbf{Descrizione} & \textbf{Stato Test} \\
    \hline
    \endfirsthead 
    \rowcolor{white}
    \caption{Tabella dei test di unità.} 
    \label{table:Tabella dei test di unità}
    \endlastfoot 
    
    TU-01 & Verificare che la classe \verb|SimlatorControllerFactory| crei correttamente un'istanza
    della classe \verb|SimulatorController|.  & V \\
    \hline

    TU-02 & Verificare che la classe \textit{SimulatorController} funzioni correttamente. Nello
    specifico, si esaminano i metodi \verb|start_all| e \verb|stop_all|.  & V \\
    \hline

    TU-03 & Verificare che la classe \verb|SimulatorThread| adoperi correttamente. & NI \\
    \hline

    TU-04 & Verificare che la classe \verb|TemperatureSensor| restituisca la stringa in formato JSON attesa. & V \\
    \hline

    TU-05 & Verificare che la classe \verb|HumiditySensor| restituisca la stringa in formato JSON attesa. & V \\
    \hline

    TU-06 & Verificare che la classe \verb|ParkingSensor| restituisca la stringa in formato JSON attesa. & NI \\
    \hline

    TU-07 & Verificare che la classe \verb|PaymentParkingSensor| restituisca la stringa in formato JSON attesa. & NI \\
    \hline
    TU-08 & Verificare che le istanze della classe \verb|KafkaProducer| adoperino nelle modalità attese nella produzione di dati, sia nel caso di success che nel caso di failure. & V \\
    \hline

    TU-09 & Verificare che la funzione \verb|write| della classe \verb|KafkaStreamWriter| funzioni correttamente. & V \\
    \hline

    TU-10 & Verificare che la funzione \verb|write| della classe
    \verb|StdoutStreamWriter| funzioni correttamente. & V \\
    \hline

    TU-11 & Verificare che la funzione \verb|acked| funzioni correttamente sia nel caso di success che nel caso di failure. & V \\
    \hline

    TU-12 & Verificare che la classe \verb|Coordinates| funzioni correttamente. & V \\
    \hline

    TU-13 & Verificare che la funzione \verb|jsonfy| funzioni correttamente. & V \\
    \hline

\end{longtable}
\clearpage
\subsection{Test di integrazione}
\renewcommand{\arraystretch}{2.5}
\rowcolors{2}{gray!20}{white}
\begin{longtable}{|>{\centering}p{2cm}|>{\RaggedRight}m{12cm}|>{\centering\arraybackslash}p{2cm}|}
    \hline
    \rowcolor{white}
    \textbf{Codice Test} & \textbf{Descrizione} & \textbf{Stato Test} \\
    \hline
    \endfirsthead 
    \rowcolor{white}
    \caption{Tabella dei test di integrazione.} 
    \label{table:Tabella dei test di integrazione}
    \endlastfoot 
        
        TI-01 & Verificare la persistenza dei dati della classe \verb|TemperatureSensor| nel \glossterm{database} \glossterm{ClickHouse}.  & V \\
        \hline
    
        TI-02 & Verificare la persistenza dei dati della classe \verb|HumiditySensor| nel \glossterm{database} \glossterm{ClickHouse}.  & V \\
        \hline

		TI-03 & Verificare la persistenza dei dati della classe \verb|ParkingSensor| nel \glossterm{database} \glossterm{ClickHouse}. & NI \\
		\hline 
		TI-04 & Verificare la persistenza dei dati della classe \verb|PaymentParkingSensor| nel \glossterm{database} \glossterm{ClickHouse}. & NI \\
    \hline
\end{longtable}
\subsection{Test di sistema}
\renewcommand{\arraystretch}{2.5}
\rowcolors{2}{gray!20}{white}
\begin{longtable}{|>{\centering}p{2cm}|>{\RaggedRight}m{12cm}|>{\centering\arraybackslash}p{2cm}|}
    \hline
    \rowcolor{white}
    \textbf{Codice Test} & \textbf{Descrizione} & \textbf{Stato Test} \\
    \hline
    \endfirsthead 
    \rowcolor{white}
    \caption{Tabella dei test di sistema.} 
    \label{table:Tabella dei test di sistema}
    \endlastfoot 
    TS-01 & Verificare che l’\glossterm{amministratore pubblico} possa accedere all’applicativo
    senza dover effettuare l’autenticazione. & NV\\
    \hline
    TS-02 & Verificare che l’\glossterm{amministratore pubblico} possa visualizzare un menù di selezione delle
    \glossterm{dashboard}, che permetta di selezionare una dashboard tra: Sensori, Ambientale e Urbanistica. & NV\\
    \hline
    TS-03 & Verificare che l’\glossterm{amministratore} pubblico possa monitorare i dati provenienti
    dai sensori relativi ai dati ambientali in una \glossterm{dashboard} apposita.
    & NV \\
    \hline
    TS-04 & Verificare che l’\glossterm{amministratore pubblico} possa visualizzare un pannello
    contenente un grafico in formato time series rappresentante l'andamento della temperatura,
    espressa in gradi Celsius (°C), per ciascun sensore, nella \glossterm{dashboard} relativa ai dati ambientali.
    & NV \\
    \hline
    TS-05 & Verificare che l’\glossterm{amministratore pubblico} possa visualizzare un pannello
    contenente un grafico in formato time series rappresentante l'andamento dell'umidità, espressa
    in percentuale, per ciascun sensore, nella \glossterm{dashboard} relativa ai dati ambientali.
    & NV \\
    \hline
    TS-06 &  Verificare che l’\glossterm{amministratore pubblico} possa visualizzare un pannello
    contenente un grafico in formato time series rappresentante l'andamento della temperatura
    percepita,
    espressa in gradi Celsius (°C), per ciascun coppia di sensori temperatura-umidità, nella
    \glossterm{dashboard} relativa ai dati ambientali.
    & NV \\
    \hline
    TS-07 & 
    Verificare che l’\glossterm{amministratore pubblico} possa visualizzare un pannello contenente
    un grafico in formato di tipo Gauge rappresentante la media aritmetica della temperatura,
    espressa in gradi Celsius (°C), per tutti i sensori, che aggreghi i dati, nella \glossterm{dashboard} relativa ai dati ambientali.
    & NV \\
    \hline 
    TS-08 & Verificare che l’\glossterm{amministratore pubblico} possa  monitorare i dati
    provenienti dai sensori relativi ai dati dei sensori in una \glossterm{dashboard} apposita.
    & NV \\
    \hline
    TS-09 & Verificare che l’\glossterm{amministratore pubblico} possa visualizzare un pannello
    contenente una mappa che mostri le posizioni dei sensori che monitorano i dati ambientali,
    mediante icone in base al tipo di sensore, nella \glossterm{dashboard} relativa ai dati
    ambientali. & NV \\
    \hline
    TS-10 & Verificare che l’\glossterm{amministratore pubblico} possa monitorare i dati provenienti
    dai sensori relativi ai dati urbanistici in una \glossterm{dashboard} apposita.
    & NV \\
    \hline
    TS-11 & Verificare che l’\glossterm{amministratore pubblico} possa visualizzare un pannello in formato status history che evidenzi lo stato dei posti nei vari parcheggi, nella
    \glossterm{dashboard} relativa ai dati urbanistici. & NV \\
    \hline
    TS-12 & 
    Verificare che l’\glossterm{amministratore pubblico} possa visualizzare un pannello in formato
    alert list che evidenzi l'elenco delle notifiche di pagamento nei vari parcheggi, nella \glossterm{dashboard} relativa ai dati urbanistici.& NV \\
    \hline
    TS-13 & Verificare che l’\glossterm{amministratore pubblico} possa filtrare i dati, visualizzati
    all’interno di un grafico di tipo time series, in base ad un sottoinsieme di sensori da lui
    selezionato. & NV \\
    \hline
    TS-14 & Verificare che l’\glossterm{amministratore pubblico} possa filtrare i dati in base ad un intervallo temporale. La \glossterm{dashboard} DI interessE deve mostrare solamente i dati aventi un timestamp in tale intervallo.
    & NV\\
    \hline
    TS-15 & Verificare che l’\glossterm{amministratore pubblico}, riceva un messaggio di errore qualora, il
    sistema di visualizzazione non riesca a reperire i dati necessari per un determinato pannello. &
    NV \\
    \hline
    TS-16 & Verificare che un sensore possa inserire le rilevazioni della temperatura, espresse in
    gradi Celsius (°C), con annesso coordinate e timestamp della rilevazione. & NV \\
    \hline
    TS-17 & Verificare che un sensore possa inserire le rilevazioni della temperatura, espresse in
    percentuale, con annesso coordinate e timestamp della rilevazione. & NV \\
    \hline 
    TS-18 & Verificare che un sensore possa inserire le rilevazioni dello stato dei parcheggi,
    espresse tramite valore booleano, con annesso coordinate e timestamp della rilevazione. & NV \\
    \hline 
    TS-19 & Verificare che un sensore possa inserire le rilevazioni dello stato dei pagamenti
    dei parcheggi con annesso coordinate e timestamp della rilevazione. & NV \\
    \hline
    TS-20 & Verificare che sia stato implementato almeno un simulatore per ogni tipologia di sensore
    & NV \\
    \hline
    TS-21 & Verificare che i dati prodotti dalle simulazioni siano realistici. & NV \\
    \hline
    TS-22 & Verificare che il sistema possa rilevare eventuali relazioni tra sorgenti di dati
    diversi. & NI \\
    \hline
    TS-23 & Verificare che il sistema possa effettuare previsioni di eventi futuri, sulla base di dati storici e attuali. & NI \\
    \hline
\end{longtable}
\subsubsection{Tracciamento dei test di sistema}
\subsection{Test di accettazione}
Questa sezione presenta i test di accettazione del software, condotti da NaN1fy e dalla
\glossterm{Proponente} sotto la supervisione del gruppo. Questi test mirano a validare il prodotto prima del suo rilascio, garantendone la \glossterm{qualità} e l'aderenza ai requisiti.

\renewcommand{\arraystretch}{2.5}
\rowcolors{2}{gray!20}{white}
\begin{longtable}{|>{\centering}p{2cm}|>{\RaggedRight}m{12cm}|>{\centering\arraybackslash}p{2cm}|}
    \hline
    \rowcolor{white}
    \textbf{Codice Test} & \textbf{Descrizione} & \textbf{Stato Test} \\
    \hline
    \endfirsthead 
    \rowcolor{white}
    \caption{Tabella dei test di accettazione.} 
    \label{table:Tabella dei test di accettazione}
    \endlastfoot  
    TA-01 & Verificare che l’\glossterm{amministratore pubblico} senza autenticazione possa:
    \begin{enumerate}
        \setlength\itemsep{0em}
        \item Usufruire dell’applicativo senza doversi iautenticare.
    \end{enumerate} & NV \\
    \hline
    TA-02 & Verificare che l’\glossterm{amministratore pubblico}, una volta entrato nell'applicativo, possa:
    \begin{enumerate}
        \setlength\itemsep{0em}
        \item Aprire il menu di selezione delle \glossterm{dashboard};
        \item Selezionare la \glossterm{dashboard} dei sensori;
        \item Visualizzare la \glossterm{dashboard};
        \item Visualizzare un pannello con una mappa che indichi, mediante icone collocate presso le coordinate di ciascun sensore, la posizione dei sensori;
        \item Visualizzare un messaggio di avvertenza di dati mancanti, all’interno del pannello, nel caso il sistema non riesca a reperire i dati.
    \end{enumerate}
    & NV \\
    \hline
    TA-03 & Verificare che l’\glossterm{amministratore pubblico}, una volta entrato nell'applicativo, possa:
    \begin{enumerate}
        \setlength\itemsep{0em}
        \item Aprire il menu di selezione delle \glossterm{dashboard};
        \item Selezionare la \glossterm{dashboard} dei dati ambientali;
        \item Visualizzare la \glossterm{dashboard};
        \item Visualizzare un pannello contenente un grafico in formato time series che mostri i
            risultati delle rilevazioni delle temperatura, espressa in gradi Celsius (°C),
            effettuate dai singoli sensori;
        \item Visualizzare un pannello contenente un grafico in formato time series che mostri i
            risultati delle rilevazioni dell’umidità, espressa in percentuale, effettuate dai
            singoli sensori;
        \item Visualizzare un pannello contenente un grafico in formato time series che mostri i
            risultati del calcolo della temperatura percepita, espressa in gradi Celsius (°C),
            effettuate dalla coppia di sensori temperatura e umidità;
        \item Visualizzare un pannello contenente un grafico di tipo Gauge che mostri il risultato
            del calcolo della temperatura media, espressa in gradi Celsius (°C), ottenuta tramite la
            media aritmetica delle rilevazioni effettuate dai singoli sensori;
        \item Visualizzare un messaggio di avvertenza di dati mancanti, all’interno del pannello, nel caso il sistema non riesca a reperire i dati.
    \end{enumerate}
    & NV \\
    \hline
    TA-04 & Verificare che l’\glossterm{amministratore pubblico}, una volta entrato nell'applicativo, possa:
    \begin{enumerate}
        \setlength\itemsep{0em}
        \item Aprire il menu di selezione delle \glossterm{dashboard};
        \item Selezionare la \glossterm{dashboard} dei dati urbanistici;
        \item Visualizzare la \glossterm{dashboard};
        \item Visualizzare un pannello in formato alert list che mostri un elenco dettagliato di
            avvisi di pagamento, includendo informazioni come la data di pagamento e il costo
            associato a ciascuna rilevazione effettuata dai sensori;
        \item Visualizzare un pannello in formato status history che mostri l'elenco degli stati dei
            parcheggi, espressi tramite valori booleani, nel tempo, rilevati dai singoli sensori;
            parcheggi, espressi tramite booleani, nel tempo effettuate dai singoli sensori;
        \item Visualizzare un messaggio di avvertenza di dati mancanti, all’interno del pannello, nel caso il sistema non riesca a reperire i dati.
    \end{enumerate}
    & NV \\
    \hline
    TA-05 & Verificare che l’\glossterm{amministratore pubblico}, una volta entrato nell’applicativo, possa:
    \begin{enumerate}
        \item Scegliere una \glossterm{dashboard} da visualizzare;
        \item Applicare dei filtri, per visualizzare solo i dati provenienti dal sottoinsieme di
            sensori selezionato, nel caso di pannelli di tipo time series; 
        \item Applicare dei filtri, per visualizzare solo i dati provenienti dal sottoinsieme di
            sensori selezionato, nel caso di pannelli di tipo alert list;
        \item Applicare dei filtri, per visualizzare solo i dati provenienti dal sottoinsieme di
            sensori selezionato, nel caso di pannelli di tipo status history;  
        \item Applicare dei filtri, per selezionare solo i dati relativi ad un definito intervallo
            di tempo, all’interno di un’intera \glossterm{dashboard}.
    \end{enumerate}
    & NV \\
    \hline
    TA-06 & Verificare che l’\glossterm{amministratore pubblico}, una volta entrato
    nell’applicativo, possa:
    \begin{enumerate}
        \item Scegliere una \glossterm{dashboard} di cui modificare il layout;
        \item Modificare il layout dei pannelli in termini di posizione di tali pannelli e della loro dimensione.
    \end{enumerate}
    & NV \\
    \hline
    TA-07 &
    Verificare che un sensore, una volta connesso al sistema, possa:
    \begin{enumerate}
        \item Inserire il risultato della rilevazione della temperatura, espressa in gradi Celsius
            (°C), con annesso il timestamp di rilevazione e le proprie coordinate geografiche.        
    \end{enumerate}
    & NV \\
    \hline
    TA-08 &
    Verificare che un sensore, una volta connesso al sistema, possa:
    \begin{enumerate}
    \item Inserire il risultato della rilevazione dell’umidità, espressa
in percentuale, con annesso il timestamp di rilevazione e le proprie coordinate geografiche.
    \end{enumerate}
    & NV \\
    \hline
    TA-09 &
    Verificare che un sensore, una volta connesso al sistema, possa:
    \begin{enumerate}
    \item  Inserire il risultato della rilevazione della presenza di auto all’interno del
        parcheggio, con annesso il timestamp di rilevazione, la targa dell'auto e le proprie coordinate geografiche. 
    \end{enumerate}
    & NV \\
    \hline
    TA-10 &
    Verificare che un sensore, una volta connesso al sistema, possa:
    \begin{enumerate}
    \item Inserire il risultato della rilevazione di avvenuto pagamento all’interno del
        parcheggio, con annesso il timestamp di rilevazione, il valore del pagamento e le proprie coordinate geografiche. 
    \end{enumerate}
    & NV \\
    \hline
\end{longtable}
\subsubsubsection{Tracciamento dei test di accettazione}
\renewcommand{\arraystretch}{2.5}
\rowcolors{2}{gray!20}{white}
\begin{longtable}{|>{\centering}p{4cm}|>{\centering\arraybackslash}p{4cm}|}
\hline
\rowcolor{white}
\textbf{Codice Test} & \textbf{Codice caso d'uso} \\
\hline
\endfirsthead
\rowcolor{white}
\caption{Tracciamento dei test di accettazione.}
\label{table:Tracciamento dei test di accettazione}
\endlastfoot
    TA-01 & UC-1 \newline
    UC-2 \newline
    UC-3 \newline
    UC-5 \newline
    UC-6 \newline
    \\
    \hline
    TA-02 & UC-0 \newline
            UC-1 \newline
            UC-1.1 \newline
            UC-4
    \\
    \hline 
    TA-03 & UC-0 \newline
            UC-2 \newline
            UC-2.1 \newline
            UC-2.2 \newline
            UC-2.3 \newline
            UC-2.4 \newline
            UC-4
    \\
    \hline
    TA-04 & UC-0 \newline
            UC-3.1 \newline
            UC-3.2
    \\
    \hline
    TA-05 & UC-5 \newline
            UC-5.1\newline
            UC-5.2
    \\
    \hline
    TA-06 & UC-6 \\
    \hline
    TA-07 & UC-7 \newline
            UC-7.1
    \\
    \hline
    TA-08 & UC-7 \newline
    UC-7.2
    \\
    \hline
    TA-09 & UC-7 \newline
    UC-7.3
    \\
    \hline
    TA-10 & UC-7 \newline
    UC-7.4
    \\
    \hline
\end{longtable}
\clearpage
\subsection{Liste di controllo}
Le liste di controllo rappresentano un prezioso strumento a disposizione del
Verificatore per l'identificazione di errori ricorrenti nella documentazione o nel codice. Integrando la descrizione del problema, facilitano la comprensione delle modifiche richieste durante la fase di revisione.
\\
Inoltre, il Verificatore ha la possibilità di aggiornare le liste di controllo nel corso del progetto, man mano che emergono nuovi errori ricorrenti.
\subsubsection{Struttura dei documenti}
\renewcommand{\arraystretch}{2.5}
\rowcolors{2}{gray!20}{white}
\begin{longtable}{|>{\centering}p{5cm}|>{\centering\arraybackslash}p{10cm}|}
\hline
\rowcolor{white}
    \textbf{Aspetto} & \textbf{Spiegazione} \\
\hline
\endfirsthead
\rowcolor{white}
\caption{Lista di controllo per la struttura dei documenti.}
\label{table:Lista di controllo per la struttura dei documenti}
\endlastfoot
    Vuoti documentativi & Non devono essere presenti sezioni senza contenuto. \\
\hline
    Didascalia assente & Tutte le tabelle e le immagini devono avere una didascalia descrittiva. \\
\hline    
Titolo principale & Tutti i titoli principali devono iniziare la pagina nella quale vengono
    inseriti \\
\hline
    Aggiornamento fantasma & Ad ogni insieme di modifichedeve corrispondere una riga nella tabella
    del changelog. \\
\hline
\end{longtable}
\subsubsection{Errori ortografici, di lingua italiana e di forma}
\renewcommand{\arraystretch}{2.5}
\rowcolors{2}{gray!20}{white}
\begin{longtable}{|>{\centering}p{5cm}|>{\centering\arraybackslash}p{10cm}|}
\hline
\rowcolor{white}
    \textbf{Aspetto} & \textbf{Spiegazione} \\
\hline
\endfirsthead
\rowcolor{white}
\caption{Lista di controllo per gli errori ortografici, di lingua italiana e di forma.}
\label{table:Lista di controllo per gli errori ortografici, di lingua italiana e di forma}
\endlastfoot
    Errori di sintassi & Gli errori di sintassi (battitura o distrazione) devono essere rimossi.\\
\hline
    Errori di coniugazione & Gli errori di coniugazione devono essere rimossi. \\
\hline
    Forma non concisa & Le espressioni troppo verbose, dove possibile, devono essere ridotte.\\
\hline
    Non formalità & Le espressioni non formali devono essere sostituite con le corrispondenti
    espressioni formali. \\
\hline
    Richiamo errato al documento & I richiami ai documenti, devono seguire la seguente forma:
    \textit{NomeDocumento vX.X.X} (e.g. \textit{Piano di Progetto v1.0.0}).\\
\hline
    Acronimi non in maiuscolo & Gli acronimi devono essere completamente in maiuscolo. \\
\hline
\end{longtable}
\subsubsection{Non conformità con le \textit{Norme di Progetto}}
\renewcommand{\arraystretch}{2.5}
\rowcolors{2}{gray!20}{white}
\begin{longtable}{|>{\centering}p{5cm}|>{\centering\arraybackslash}p{10cm}|}
\hline
\rowcolor{white}
    \textbf{Aspetto} & \textbf{Spiegazione} \\
\hline
\endfirsthead
\rowcolor{white}
    \caption{Lista di controllo per le non conformità con le \textit {Norme di Progetto}.}
    \label{table: Lista di controllo per le non conformità con le Norme di Progetto}
\endlastfoot
    Formato date errato & Il formato delle date deve essere aaaa-mm-dd all’interno dei documenti. \\
    \hline
    Punteggiatura scorretta negli elenchi &  Ogni elemento di un elenco, numerato o non, deve ter-
    minare con un ``;”, ad eccezione dell’ultima riga, la quale deve terminare con ``.”. \\ 
    \hline
    ``:” non in grassetto negli elenchi & Gli elenchi nella forma ``termine: testo”, devono
    includere ``:” nel grassetto. \\
    \hline
    Maiuscole nei titoli & La prima lettera di ogni titolo deve essere maiuscola. Il resto del
    titolo dovrebbe essere in minuscolo (tolte particolari eccezioni). \\
    \hline
    Ruoli in minuscolo & Tutti i ruoli del progetto devono avere la prima lettera in maiuscolo. \\
    \hline
    Termine non presente nel glossario & Ogni termine segnato con la formattazione da glossario deve essere presente nel glossario. \\
\hline
\end{longtable}
\newpage
\section{Cruscotto della qualità}
\subsection{Qualità di processo - Fornitura}
\subsubsection{MPC-EAC Estimate at Completion}
\textbf{\glossterm{RTB}:} EAC indica, per ogni periodo, la stima del budget finale prefiguratasi rispetto al valore prefissato definito dal BAC ovvero il Budget at Completion. Questo valore viene calcolato come il rapporto tra BAC e CPI dove CPI è Cost Performance Index.
\subsubsection{MPC-EV Earned Value e MPC-PV Planned Value}
\textbf{\glossterm{RTB}:} EV rappresenta il valore generato dal prodotto sviluppato fino a un dato momento mentre PV costituisce il valore del lavoro pianificato fino a un dato momento. Il loro confronto permette di comprendere lo scostamento che avviene tra i preventivi e i consuntivi di lavoro effettuati e quindi ottenere una chiara visione dell'andamento del lavoro svolto dal team. 
\subsubsection{MPC-BV Budget Variance e MPC-SV Schedule Variance}
\textbf{\glossterm{RTB}:} BV consiste nella differenza tra il valore guadagnato (EV) e i costi effettivamente sostenuti (AC). Un valore negativo di BV indica che si sta spendendo più di quanto si stia guadagnando. SV, invece, rappresenta la differenza tra il valore guadagnato (EV) e il valore pianificato (PV). Un valore negativo indica un ritardo rispetto alla produzione e consegna del lavoro preventivato.
\subsubsection{MPC-AC Actual Cost e MPC-ETC Estimate To Complete}
\textbf{\glossterm{RTB}:} AC rappresenta il costo sostenuto fino a un dato momento mentre ETC consiste nella conseguente stima dei costi da sostenere per il completamento del progetto.
\subsubsection{MPC-CPI Cost Performance Index}
\textbf{\glossterm{RTB}:} CPI indica il rapporto tra il valore guadagnato (EV) e i costi sostenuti (AC). Un valore di CPI uguale a 1 indica che i costi sostenuti sono in linea con il bugdet pianificato (BAC) mentre un valore negativo comporta il discostamento della stima del bidget al completamento (EAC) dal budget effettivo. 
\subsection{Qualità di processo - Documentazione}
\subsubsection{Indice Gulpease}
\textbf{\glossterm{RTB}:}
\subsubsection{Correttezza Ortografica}
\textbf{\glossterm{RTB}:}
\subsection{Qualità di processo - Gestione della qualità}
\subsubsection{MPC-QMS Quality Metrics Satisfied}
\textbf{\glossterm{RTB}:}
\subsection{Qualità di processo - Gestione dei processi}
\subsubsection{MPC-NR Non-calculated Risk}
\textbf{\glossterm{RTB}:}
\subsection{Qualità di prodotto - Funzionalità}
\textbf{\glossterm{RTB}:}
\subsection{Qualità di prodotto - Affidabilità}
\textbf{\glossterm{RTB}:}
\subsection{Qualità di prodotto - Usabilità}
\textbf{\glossterm{RTB}:}
\subsection{Qualità di prodotto - Efficacia}
\textbf{\glossterm{RTB}:}
\subsection{Qualità di prodotto - Manutenibilità}
\textbf{\glossterm{RTB}:}
\end{document}
